
\subsubsection{6.1 Interpretation of Infrastructure Findings}

\textbf{The EvalPlus failure set constitutes a reproducible foundation for repair research.} The four-condition verification protocol confirms that all components required for mechanism experiment execution --- failure identifiers, stored incorrect solutions, EvalPlus API access, and deterministic test selection --- are available from the h-e1 Run 2 archive. Any research group with access to the archive and EvalPlus 0.3.1 can reproduce the infrastructure without additional baseline API calls.

\textbf{Static analysis is the wrong oracle for EvalPlus semantic failures.} The 16--32\% fire rate across both benchmark subsets means that the majority of failures are opaque to the standard tool-assisted repair pipeline. This falsification of the static analysis oracle, established empirically in h-e1 Run 2, is the primary motivation for specification-aligned repair as a replacement oracle. Specification-aligned repair addresses the failure mode directly: it provides the intent, behavioral gap, and deviation signal that static analysis cannot supply.

\textbf{The h-e1 to h-e1-v2 trajectory illustrates appropriate verification target selection.} The initial h-e1 verification script tested \texttt{plus\textbackslash \_fail\textbackslash \_tests} completeness and failed because 6/134 records had empty fields (gate: FAIL, 5/6 checks passing). The h-e1-v2 redesign correctly identified that \texttt{solutions\textbackslash \_cache.jsonl} --- not \texttt{plus\textbackslash \_fail\textbackslash \_tests} --- is the relevant data source for prompt construction in Conditions B and C, because the model's incorrect output is stored in the solutions cache and the failing test input/output pair is obtained via the EvalPlus API at experiment time. This redesign represents appropriate hypothesis refinement in response to a verification failure, not abandonment of the core research direction.

\subsubsection{6.2 Limitations}

\textbf{Mechanism experiments not executed.} Predictions P1, P2, and P3 are pre-registered but unmeasured. All comparative claims --- whether specification-aligned repair outperforms blind reprompting --- are INCONCLUSIVE. This paper establishes infrastructure and registers the design; the mechanism experiments are the primary next step.

\textbf{Six-task exclusion (n = 128, not n = 134).} The principled exclusion of 4.5\% of the failure set for missing archive data reduces the working set modestly. Statistical power analysis indicates that n = 128 is adequate for the expected effect size range (15--40\% fix rate under Condition C). The exclusion is documented and the specific task identifiers are reported.

\textbf{Single model and temperature.} All experiments use GPT-4o-mini at temperature = 0.2, seed = 42. Generalization to other models, model scales, temperatures, and sampling configurations is not addressed and is left for future work. The assumption that GPT-4o-mini can leverage structured specification context to improve repair performance over blind reprompting remains unverified pending mechanism experiment execution.

\textbf{Component ablation deferred.} The specification triple is evaluated as an integrated oracle. Individual component attribution --- whether docstring re-anchoring, input/output counterexample, or actual output deviation detection is the primary driver of any repair improvement --- requires a full 7-condition ablation. This ablation is deferred to future work; the present design tests only the integrated triple against a blind reprompting placebo.

\textbf{Single-round repair.} Only round-1 repair is evaluated under conditions B and C. Multi-round repair with updated specification context per round is a tractable extension deferred to future work.

\subsubsection{6.3 Implications for Reproducibility Methodology}

The four-condition existence verification protocol developed for this work --- checking failure identifier availability, stored solution coverage, API accessibility, and deterministic test selection --- is a reusable template applicable to any repair experiment built on an archived LLM evaluation failure set. Pre-registration of quantitative predictions before mechanism data collection prevents hypothesis drift and guards against post-hoc analysis bias. Together, the verified infrastructure and the pre-registered design constitute a methodological contribution to reproducibility in LLM code repair evaluation that is independent of the mechanism experiment outcomes.

